% Created 2026-10-15 gio 22:53
% Intended LaTeX compiler: pdflatex
\documentclass[11pt]{article}
\usepackage[utf8]{inputenc}
\usepackage[T1]{fontenc}
\usepackage{graphicx}
\usepackage{longtable}
\usepackage{wrapfig}
\usepackage{rotating}
\usepackage[normalem]{ulem}
\usepackage{amsmath}
\usepackage{amssymb}
\usepackage{capt-of}
\usepackage{hyperref}
\usepackage[newfloat]{minted}
\author{marco robba}
\date{\today}
\title{Lezioni: errori di design e di architettura}
\hypersetup{
 pdfauthor={marco robba},
 pdftitle={Lezioni: errori di design e di architettura},
 pdfkeywords={},
 pdfsubject={},
 pdfcreator={Emacs 30.2 (Org mode 9.7.29)}, 
 pdflang={Italian}}
\begin{document}

\maketitle
\tableofcontents

Gli errori di design e di architettura fatti durante lo sviluppo del gioco, in ordine, ciascuno con la soluzione che
ha funzionato. Prova: l'hash del commit, oppure il documento dove l'argomento è trattato per intero. Le lezioni di
processo (agenti, merge, build) sono in \href{WORKFLOW.md}{WORKFLOW.md} e \href{DEV\_GUIDE.md}{DEV\textsubscript{GUIDE.md}}; i dettagli sulla
taratura dell'IA in \href{AI\_TUNING.md}{AI\textsubscript{TUNING.md}}. L'ultimo capitolo, \hyperref[sec:org0c05727]{Costi}, misura quanti token è costato ogni modulo
e ogni scopo.

\noindent\rule{\textwidth}{0.5pt}
\section{2026-10-10}
\label{sec:org8b5fb2f}

\subsection{Antiaerea contro aerei che girano in cerchio}
\label{sec:org988693b}
\begin{itemize}
\item Errore: antiaerea e missili prevedevano un volo in linea retta, quindi gli aerei che orbitavano attorno a un
bersaglio non venivano mai colpiti.
\item Soluzione: previsione d'anticipo che tiene conto della virata, missili a inseguimento con anticipo, spoletta della
contraerea più intelligente (\texttt{d8eea2e}).
\end{itemize}
\section{2026-10-12 → 14: IA delle schermaglie}
\label{sec:org0d22547}

\subsection{Partite che non finivano mai}
\label{sec:org96e69a5}
\begin{itemize}
\item Errore: i comandanti restavano a casa dietro le difese, e eserciti alla pari si scambiavano perdite uguali sul
campo all'infinito.
\item Soluzione: una partita finisce solo quando muore un comandante, quindi a fine partita l'IA spinge verso l'ultima
posizione nota del comandante nemico e lo elimina. Ora finisce l'86 \% delle partite Normal contro Normal.
\end{itemize}
\subsection{Raggruppamento dopo una sortita}
\label{sec:orgd40860e}
\begin{itemize}
\item Errore: dopo una sortita l'esercito tornava al suo punto di raccolta. Finivano meno partite.
\item Soluzione: modifica annullata. È la pressione continua, a ondate sparse, che fa finire le partite su Highlands e
Canyon (\texttt{44ba888}).
\end{itemize}
\subsection{Composizione dell'esercito di Hard}
\label{sec:orgbbb9024}
\begin{itemize}
\item Errore: troppa poca artiglieria nell'esercito di Hard (\texttt{ArtilleryShare 0.15}: 69 vittorie, 105 sconfitte).
\item Soluzione: \texttt{ArtilleryShare 0.55} (166 vittorie, 8 sconfitte, 18 pareggi su 192). Dipende dal fatto che
l'artiglieria T1 rende moltissimo rispetto al suo costo, quindi va ritarato dopo un intervento sul bilanciamento.
\end{itemize}
\section{2026-10-13}
\label{sec:orgaa2a005}

\subsection{Attaccare bersagli persi nella nebbia}
\label{sec:orgc09c19b}
\begin{itemize}
\item Errore: le unità mantenevano l'ordine di attacco su un bersaglio che non vedevano più.
\item Soluzione: vanno all'ultima posizione nota, poi abbandonano l'ordine (\texttt{7c232bd}).
\end{itemize}
\subsection{Indici dei blueprint}
\label{sec:orgdcd6558}
\begin{itemize}
\item Errore: il codice usava gli indici dei blueprint. I blueprint si caricano in ordine di nome file, quindi aggiungere
\texttt{uef\_t1\_transport} ha spostato tutte le unità successive.
\item Soluzione: cercare i blueprint per id, oppure per ruolo di fazione (vedi sotto).
\end{itemize}
\subsection{Esplosioni alla morte}
\label{sec:orgcdb08df}
\begin{itemize}
\item Errore: l'esplosione alla morte di un'unità sperimentale veniva aggiunta alla lista dei danni da \texttt{OnDeath}, mentre
il sistema dei danni stava scorrendo quella stessa lista.
\item Soluzione: mettere l'esplosione in coda ed emetterla in un tick successivo (\texttt{73d9b54}).
\end{itemize}
\subsection{Grandi unità mobili}
\label{sec:org1fad928}
\begin{itemize}
\item Errore: il raggio mobile più grande stabilisce la dimensione della ricerca di separazione per tutte le unità,
quindi una sola unità grande rallentava l'intera simulazione.
\item Soluzione: \texttt{Radius} delle unità mobili ≤ 3.0; le unità grandi hanno un raggio di contatto separato
(\texttt{CrushRadius}).
\end{itemize}
\subsection{Collisioni tra navi}
\label{sec:org891c4f9}
\begin{itemize}
\item Errore: le navi si sovrapponevano tra loro.
\item Soluzione: separazione a capsula di scafo tra le navi di tutti gli eserciti; i sottomarini collidono solo con i
sottomarini (\texttt{cbf7f9a}).
\end{itemize}
\subsection{Stato della simulazione salvabile}
\label{sec:org5f492bd}
\begin{itemize}
\item Errore: delegate in cache e liste temporanee statiche nello stato della simulazione non si potevano salvare.
\item Soluzione: sono marcati \texttt{[NoSave]} e ricostruiti quando servono; \texttt{SaveCoverageTests} indica per nome ogni campo
che non si può salvare (\texttt{da42cb8}). Regola: CLAUDE.md, regola rigida 9.
\end{itemize}
\section{2026-10-14}
\label{sec:orge613f85}

\subsection{Id delle unità scritti nel codice}
\label{sec:orgb11678c}
\begin{itemize}
\item Errore: il codice della simulazione e dell'IA nominava direttamente le unità UEF (\texttt{uef\_...}), quindi non si
poteva aggiungere una seconda fazione.
\item Soluzione: ruoli di fazione: \texttt{RoleBlueprint(faction, role)}, \texttt{-1} = la fazione non ne ha (\texttt{20963ba}).
Contratto: ARCHITECTURE ``Factions''.
\end{itemize}
\subsection{Conflitti tra tasti rapidi}
\label{sec:org22850f6}
\begin{itemize}
\item Errore: due funzioni usavano gli stessi tasti (D, E: immersione e cattura contro attracco e trasporto).
\item Soluzione: tasti rapidi dipendenti dal contesto, decisi in base alla selezione (\texttt{Input/SharedHotkeys.cs},
\texttt{e747ee9}).
\end{itemize}
\section{Costi}
\label{sec:org0c05727}

Quanti token è costato costruire il gioco fin qui: per modulo, per scopo e per tipo di lavoro, e che cosa insegnano
questi numeri. Dati dal 2026-10-10 al 2026-10-15 alle 22:30; l'onda 5 è ancora in corso, quindi i suoi numeri sono
parziali.
\subsection{Come sono stati misurati}
\label{sec:org81363e2}
\begin{itemize}
\item Fonte: le trascrizioni di Claude Code di questo progetto (\texttt{\textasciitilde{}/.claude/projects/-run-media-marco-Volume-D-RTS/}):
14 sessioni principali (orchestratore) e 81 agenti (sub-agenti e sub-sub-agenti), 6.003 chiamate al modello.
Ogni chiamata è contata una volta sola. Ogni agente è assegnato a un modulo in base alla descrizione con cui è
stato lanciato; la sessione principale conta come ``Fondazione'' fino al lancio dei primi agenti (2026-10-10 19:24),
poi come ``Orchestrazione''.
\item Token di input (scritture in cache, letture dalla cache, input normale): esatti, presi dalle trascrizioni.
\item Token di output: \textbf{stimati}. In circa 4 chiamate su 5 la trascrizione registra solo il conteggio parziale
dell'inizio dello stream. Per quelle chiamate l'output è stimato dal testo e dalle chiamate agli strumenti visibili
(caratteri / 4) × 2,53, il rapporto misurato sulle 1.189 chiamate con il conteggio completo (comprende il
ragionamento nascosto). Errore plausibile: ±30 \% sull'output, circa ±7 \% sul costo totale.
\item Dollari: equivalente a listino API di Claude Opus 5.5 (per milione di token: input 4 \$, scrittura in cache 5 \$
(5 minuti) o 8 \$ (1 ora), lettura dalla cache 0,20 \$, output 20 \$). Il proprietario usa un abbonamento e non paga
questa cifra: i dollari servono a confrontare i moduli tra loro. Il vincolo reale dell'abbonamento sono i limiti
di sessione (vedi \hyperref[sec:org196e259]{Costo del processo}).
\item Righe: righe aggiunte in \texttt{src/}, \texttt{tests/}, \texttt{tools/}, \texttt{data/} dal merge del modulo in \texttt{main} (diff rispetto al
primo genitore). Per l'onda 5, solo le righe già committate sul ramo dell'agente.
\end{itemize}
\subsection{Totale}
\label{sec:org5347403}
In 6 giorni: \textbf{821,6 milioni di token, circa 423 \$} (in media 70 \$ al giorno), per circa 65.000 righe.

\begin{center}
\begin{tabular}{lllrr}
Componente & Token (M) & \% token & \$ & \% costo\\
\hline
Lettura dalla cache & 784,4 & 95,5 & 157 & 37\\
Scrittura in cache & 32,5 & 4,0 & 174 & 41\\
Output (stimato) & 4,6 & 0,6 & 92 & 22\\
\end{tabular}
\end{center}

L'input fuori cache è trascurabile (12.000 token).

\begin{center}
\begin{tabular}{rlrr}
Giorno & Token (M) & \$ & Chiamate\\
\hline
10-10 & 89,8 & 52 & 483\\
10-11 & 192,0 & 84 & 875\\
10-12 & 204,4 & 105 & 1.673\\
10-13 & 142,4 & 82 & 1.184\\
10-14 & 91,7 & 46 & 843\\
10-15 & 101,3 & 54 & 945\\
\end{tabular}
\end{center}

Il picco del 12 ottobre è l'onda 3: cinque agenti in parallelo, ciascuno con i propri sub-agenti, più le loro
riprese.
\subsection{Per modulo}
\label{sec:org8766aed}
Ordinati per costo. ``\$ / 1000 righe'' confronta l'impegno con il risultato; per i moduli dell'onda 5 non è ancora
significativo.

\begin{center}
\begin{tabular}{lrlrllrl}
Modulo & Agenti & Token (M) & Output (k) & \$ & \% costo & Righe & \$ / 1000 righe\\
\hline
Orchestrazione (sessioni principali) & — & 81,0 & 518 & 56,0 & 13,2 & — & —\\
Rifinitura SupCom (4f) & 9 & 79,5 & 404 & 41,3 & 9,8 & 3.793 & 10,9\\
IA schermaglie (onde 2-4) & 7 & 90,5 & 389 & 39,8 & 9,4 & 4.621 & 8,6\\
Navale (3c) & 4 & 61,0 & 229 & 31,5 & 7,4 & 3.143 & 10,0\\
Mappe, menu, flusso (onda 2) & 3 & 64,8 & 231 & 24,6 & 5,8 & 3.807 & 6,5\\
Armi strategiche (3b) & 6 & 45,7 & 254 & 22,3 & 5,3 & 4.343 & 5,1\\
Sperimentali e comandanti (4g) & 5 & 38,3 & 232 & 22,1 & 5,2 & 3.006 & 7,3\\
Trasporti (4h) & 4 & 37,9 & 225 & 20,1 & 4,8 & 3.350 & 6,0\\
Salvataggi e replay (4i) & 3 & 30,9 & 210 & 16,8 & 4,0 & 4.473 & 3,8\\
Sandbox (3e) & 5 & 28,1 & 211 & 14,4 & 3,4 & 4.025 & 3,6\\
UX (5n, in corso) & 6 & 30,6 & 171 & 14,3 & 3,4 & 1.753 & —\\
Scala e prestazioni (5m, in corso) & 5 & 26,0 & 139 & 13,3 & 3,2 & 810 & —\\
Nebbia di guerra e intel (onda 2) & 2 & 30,9 & 135 & 12,7 & 3,0 & 1.332 & 9,5\\
Fazione Cybran (5l, in corso) & 5 & 22,5 & 106 & 11,1 & 2,6 & 68 & —\\
Tecnologia T2/T3 (3a) & 2 & 22,2 & 120 & 10,5 & 2,5 & 3.125 & 3,3\\
Grafica client (onda 1) & 1 & 20,7 & 162 & 9,9 & 2,3 & 3.361 & 2,9\\
Audio e rifinitura (onda 2) & 2 & 20,6 & 107 & 9,7 & 2,3 & 2.622 & 3,7\\
Pathfinding e movimento (onda 1) & 1 & 18,1 & 128 & 8,9 & 2,1 & 2.428 & 3,7\\
Multigiocatore lockstep (5o, in corso) & 4 & 13,8 & 126 & 8,6 & 2,0 & 2.854 & —\\
Combattimento (onda 1) & 1 & 11,9 & 66 & 7,3 & 1,7 & 1.506 & 4,8\\
UI e HUD (onda 1) & 1 & 12,0 & 97 & 6,6 & 1,6 & 2.674 & 2,5\\
IA usa tutto (5k, in corso) & 2 & 12,9 & 51 & 6,2 & 1,5 & 368 & —\\
Economia e costruzione (onda 1) & 1 & 6,7 & 77 & 5,5 & 1,3 & 1.644 & 3,3\\
Infrastruttura fazioni (j) & 2 & 11,0 & 81 & 5,1 & 1,2 & 921 & 5,6\\
Fondazione (sessione principale) & — & 3,9 & 156 & 5,1 & 1,2 & 5.113 & 1,0\\
\hline
\textbf{Totale} & 81 & 821,6 & 4.625 & 423,5 & 100,0 & 65.140 & —\\
\end{tabular}
\end{center}
\subsection{Per scopo}
\label{sec:orgd1a423c}
Gli stessi moduli raggruppati per lo scopo che servono.

\begin{center}
\begin{tabular}{lrrrrl}
Scopo & Token (M) & \$ & \% costo & Righe & \$ / 1000 righe\\
\hline
Gameplay SupCom (unità e sistemi) & 386 & 198 & 47 & 28.659 & 6,9\\
Esperienza client (grafica, UI, audio, menu, sandbox) & 177 & 79 & 19 & 18.242 & 4,4\\
Coordinamento (orchestratore) & 81 & 56 & 13 & — & —\\
Avversario IA & 103 & 46 & 11 & 4.989 & 9,2\\
Infrastruttura tecnica (fondazione, salvataggi, rete, prestazioni) & 75 & 44 & 10 & 13.250 & 3,3\\
\end{tabular}
\end{center}

Gameplay: economia, combattimento, pathfinding, nebbia, T2/T3, strategiche, navale, sperimentali, trasporti,
rifinitura 4f, fazioni, Cybran. Client: grafica, UI e HUD, audio, sandbox, UX, mappe e menu. Infrastruttura:
fondazione, salvataggi e replay, multigiocatore, scala e prestazioni.

Metà della spesa è andata al gameplay. L'avversario IA è lo scopo più caro per riga (9,2 \$ ogni 1000 righe),
l'infrastruttura il più economico (3,3 \$): è quasi tutto codice nuovo e separato dal resto.
\subsection{Costo del processo}
\label{sec:org196e259}
Quanto è costato il lavoro che non scrive direttamente il gioco.

\begin{center}
\begin{tabular}{lllrr}
Voce & Quanti & Token (M) & \$ & \% costo\\
\hline
Agenti al primo lancio & 58 agenti & 523,6 & 263 & 62\\
Riprese (fresh, resume, finish, part 2/3, f2/g2/h2/i2/j2) & 20 agenti & 179,0 & 86 & 20\\
Sessioni principali (orchestratore, fondazione) & 14 sessioni & 84,8 & 61 & 14\\
Agenti di solo merge & 3 agenti & 34,3 & 13 & 3\\
\end{tabular}
\end{center}

\begin{itemize}
\item Avvio di un agente: in mediana un agente fa 24 chiamate, circa 1 \$, prima della sua prima modifica a un file
(lettura di CLAUDE.md, dei docs e del codice). Sono circa 79 \$ in tutto, il 24 \% del costo degli agenti.
\item Sub-sub-agenti (lanciati da un agente): 90 \$, il 21 \% del totale.
\item Limiti dell'abbonamento: le trascrizioni registrano almeno 15 limiti di sessione e 1 limite settimanale. Hanno
interrotto 76 trascrizioni tra sessioni e agenti. Il limite vale per tutto l'account, quindi tutti gli agenti in
parallelo si fermano insieme, e quasi tutte le riprese della tabella vengono da qui o dal tetto di 300k token di
contesto.
\end{itemize}
\subsection{Cosa insegnano i numeri}
\label{sec:org39dae0d}

\subsubsection{Il contesto costa più del codice}
\label{sec:org4e11c83}
\begin{itemize}
\item Dato: il 95,5 \% dei token sono riletture dalla cache; l'output, cioè il codice e il testo scritti, è il 22 \% del
costo. Ogni chiamata rilegge tutto il contesto dell'agente, quindi il costo segue la dimensione del contesto per il
numero di turni, non le righe scritte.
\item Lezione: contesti piccoli e pochi turni costano meno. Il tetto di 300k token per agente (WORKFLOW) serve anche
a contenere il costo, non solo la qualità.
\end{itemize}
\subsubsection{Il codice nuovo costa meno del codice che tocca sistemi esistenti}
\label{sec:org4ffb9c3}
\begin{itemize}
\item Dato: fondazione 1,0 \$ ogni 1000 righe, onda 1 da 2,5 a 4,8 \$. Funzioni che toccano molti sistemi esistenti
costano da 2 a 4 volte tanto: rifinitura 4f 10,9 \$, navale 10,0 \$, nebbia 9,5 \$, IA 8,6 \$.
\item Lezione: il costo per riga cresce con il codebase, perché ogni agente deve prima leggere i sistemi che tocca. Un
brief che indica già i file e i contratti (ARCHITECTURE) fa risparmiare lettura.
\end{itemize}
\subsubsection{Ogni agente in più ha un costo fisso}
\label{sec:orgd8920ce}
\begin{itemize}
\item Dato: circa 1 \$ di avvio per agente. La rifinitura 4f ha usato 9 agenti (sub-agenti, ripresa f2, merge) per
3.793 righe ed è il modulo più caro per riga.
\item Lezione: spezzare un lavoro in sub-agenti conviene solo se ogni pezzo è grande e indipendente. I lavori piccoli
conviene farli fare all'agente che ha già il contesto.
\end{itemize}
\subsubsection{Le riprese costano un quinto del totale}
\label{sec:org58e6128}
\begin{itemize}
\item Dato: 20 agenti di ripresa, 86 \$ (20 \%). Ogni ripresa riparte da zero e rilegge handoff, docs e codice.
\item Lezione: brief dimensionati per finire dentro una finestra di sessione; handoff scritti prima del limite (regola
già in WORKFLOW); non lanciare un'onda intera poco prima di un limite di sessione.
\end{itemize}
\subsubsection{L'IA costa in turni, non in righe}
\label{sec:orgb9278db}
\begin{itemize}
\item Dato: l'IA è il modulo con più token (90,5 M) per 4.621 righe. La taratura gioca centinaia di partite headless e
ogni giro misura → modifica → misura rilegge il contesto.
\item Lezione: per l'IA conta tenere corto quello che ogni giro rilegge: dagli strumenti di misura basta il riassunto
(vittorie, sconfitte, durata, costo per tick), non i log.
\end{itemize}
\subsubsection{Coordinamento}
\label{sec:orgccc3c65}
\begin{itemize}
\item Dato: orchestratore 56 \$ (13 \%), agenti di solo merge 13 \$ (3 \%): il 16 \% del totale per coordinare 81 agenti.
\item Lezione: è un costo fisso del lavoro a onde. Va tenuto basso con report brevi degli agenti e docs che
l'orchestratore non deve rileggere ogni volta.
\end{itemize}
\subsubsection{Onda 5 (in corso)}
\label{sec:org879ea6f}
\begin{itemize}
\item Dato: 53,5 \$ finora per k, l, m, n, o. Gran parte del lavoro è ancora non committato nei worktree, fermato dal
limite di sessione delle 15:00, quindi il costo per riga non è ancora significativo (per esempio Cybran: 11 \$ per
68 righe committate).
\item Lezione: rifare il conto dopo il merge dell'onda 5.
\end{itemize}
\end{document}
